\documentclass[11pt,a4paper]{article}

\usepackage[margin=2.4cm]{geometry}
\usepackage[T1]{fontenc}
\usepackage{amsmath,amssymb}
\usepackage{graphicx}
\usepackage{booktabs}
\usepackage{siunitx}
\usepackage{caption}
\usepackage{subcaption}
\usepackage{xcolor}
\usepackage[colorlinks=true,linkcolor=blue!50!black,citecolor=blue!50!black,urlcolor=blue!50!black]{hyperref}
\usepackage{listings}
\usepackage{enumitem}
\setlist{itemsep=1pt,topsep=2pt}

\lstset{
  basicstyle=\ttfamily\footnotesize,
  breaklines=true,
  frame=single,
  language=Python,
  showstringspaces=false,
}

\title{Operations Optimisation Assignment\\
\large Reproducing the MILP for the Drone-Assisted Pickup and Delivery Problem\\
{\small Mulumba \& Diabat, \emph{Transp. Res. E} 181 (2024) 103377}}
\author{Kai\\TU Delft Aerospace Engineering}
\date{May 2026}

\newcommand{\R}{\mathbb{R}}
\newcommand{\Nplus}{N^{+}}
\newcommand{\NO}{N^{0}}
\newcommand{\Cprime}{C'}
\newcommand{\E}{\mathcal{E}}
\newcommand{\dapdp}{\textsc{DAPDP}}
\newcommand{\pdp}{\textsc{PDP}}

\begin{document}
\maketitle

\section{Introduction}

This report reproduces the mixed-integer linear program of Mulumba and
Diabat~\cite{mulumba2024} for the Drone-Assisted Pickup and Delivery Problem
(\dapdp). A truck performs a sequence of paired pickup and delivery requests
while a single on-board drone may peel off, deliver a light parcel, and
rendezvous with the truck downstream. The goal is to minimise total fuel cost
subject to truck capacity, drone payload and endurance, and a maximum route
duration.

The contributions of this work are:
\begin{enumerate}
\item A clean Python/Gurobi implementation of all $37$ constraints of the
paper (Section~3.2 of the original) plus the optional valid inequalities
(Section~3.3).
\item A truck-only baseline obtained by restricting the same model with
$\Cprime{=}\emptyset$, used to compute the saving $\mathit{SPDP} = 1 -
z_{\dapdp} / z_{\pdp}$.
\item A verification suite of seven independently-audited test scenarios
covering each major class of constraint.
\item A sensitivity analysis on drone endurance, drone speed and the drone
cost factor $\alpha$.
\item A small-scale validation that recovers the same order of saving (mean
SPDP $\approx$ 18--20\%) reported in the paper.
\end{enumerate}

The paper contains a small number of typographic ambiguities
(constraint~(28), the demand sign convention in \S5.2, and the launch-side
treatment in constraints~(10) and~(26)) which we resolve and document
explicitly in \S\ref{sec:interpretations}.

\section{Problem statement}

Let $\Nplus = \{1,\dots,2n+1\}$, $\NO = \{0,\dots,2n\}$, $C =
\{1,\dots,2n\}$ be the customer node set, partitioned into pickups
$\Phi^{+} = \{1,\dots,n\}$ and deliveries $\Phi^{-} = \{n{+}1,\dots,2n\}$;
delivery node $n{+}i$ is paired with pickup $i$. Node $0$ is the start
depot, node $2n{+}1$ a virtual end depot at the same location. The
drone-eligible subset $\Cprime \subseteq \Phi^{-}$ contains those
deliveries whose demand $\le Q'$, the drone payload limit. Truck travel
time is $\tau_{ij} = d_{ij}/s$, drone travel time $\tau'_{ij} = d_{ij}/s'$
with $d_{ij}$ the Euclidean distance.

\paragraph{Demand sign convention.}
We use $q_i > 0$ at pickups and $q_i < 0$ at deliveries. This is the
\emph{opposite} of the textual statement in the paper's \S5.2 but is the
only convention consistent with the literal reading of the weight
balance~(28), where positive contributions of $q_j$ must increase the
truck load on arrival.

\paragraph{Decision variables.}
\begin{itemize}
\item $x^{v}_{ij}\in\{0,1\}$: truck $v$ traverses arc $(i,j)$.
\item $y^{v}_{ijk}\in\{0,1\}$: a drone-sortie launch $i$, drone delivery
$j$, rendezvous $k$ for truck $v$.
\item $p^{v}_{ij}\in\{0,1\}$: $i$ precedes $j$ in truck $v$'s route.
\item $t^{v}_{i},\;t'^{\,v}_{i}\ge 0$: truck and drone arrival times at $i$.
\item $u^{v}_{i}\in\{1,\dots,|N|\}$: position of $i$ in truck $v$'s
sequence.
\item $w^{v}_{i}\ge 0$: truck weight after visiting $i$.
\item $\eta\ge 0$: epigraph variable; objective $\min \eta$ with $\eta \ge
\sum_{v,i,j} c_{ij}x^{v}_{ij} + \sum_{v,(i,j,k)\in P}
(c'_{ij}{+}c'_{jk}) y^{v}_{ijk}$.
\end{itemize}

The model's $37$ constraints are stated in the paper, equations~(2)--(37); a
brief catalogue is given in Table~\ref{tab:constraints}. We refer the reader
to the paper for the full algebra and only highlight the non-trivial
implementation decisions in the next section.

\begin{table}[t]
\centering
\caption{Constraint groups of the \dapdp{} MILP and their implementation.}
\label{tab:constraints}
\small
\begin{tabular}{@{}lll@{}}
\toprule
Eqs.        & Purpose                                          & Implementation note \\
\midrule
(2)         & Each customer served exactly once                & --- \\
(3)--(4)    & Depot start/end at most once per truck           & --- \\
(5)         & Sub-tour elimination (MTZ on $u$)                & big-$M_{\text{pos}} = 2n+2$ \\
(6)         & Truck flow conservation                          & --- \\
(7)--(8)    & At most one launch/return per node               & --- \\
(9)--(10)   & Launch and rendezvous on truck route             & (10) subsumed by (9), see \S\ref{sec:interpretations} \\
(11)        & Launch precedes rendezvous in $u$                & --- \\
(12)--(15)  & Time synchronisation of truck and drone          & big-$M_t = T_{\max}+10$\,h \\
(16)        & Truck-arrival recursion (with launch/recovery)   & --- \\
(17)--(18)  & Drone-arrival recursion                          & --- \\
(19)        & Drone endurance $\le e$                          & also pre-filtered into $P$ \\
(20)--(22)  & Consistency of $u$ and $p$                       & --- \\
(23)        & No second sortie before return of first          & --- \\
(24)        & Forbid empty depot-to-depot trip                 & --- \\
(25)        & Same vehicle for paired (i, $i+n$)               & --- \\
(26)        & Pickup before launch when drone delivers $j$     & extended to $i{=}0$, see \S\ref{sec:interpretations} \\
(27)        & Pickup precedes delivery on truck                & --- \\
(28)        & Truck weight evolution                           & sign correction, see \S\ref{sec:interpretations} \\
(31)--(34)  & Time bounds, $T_{\max}$                          & --- \\
(35)        & $p^{v}_{0j}=1$                                   & omitted; see \S\ref{sec:interpretations} \\
(36)--(37)  & Bounds on $u, w$                                 & encoded in variable bounds \\
(43)--(44)  & Valid inequalities (Section 3.3)                 & enabled by default \\
\bottomrule
\end{tabular}
\end{table}

\section{Interpretation decisions where the paper is ambiguous}
\label{sec:interpretations}

The paper has a few places where the printed model is either
mathematically inconsistent or requires assumptions that are not stated.
Implementing it verbatim would yield an infeasible or trivially-relaxed
program, so we deviate as follows.

\paragraph{(28) Truck weight balance --- big-$M$ sign.}
The paper writes
\[
w^{v}_{j} \ge w^{v}_{i} + q_{j} + \sum_{(i,j,k)\in P} q_{m}\,y^{v}_{j m k}
\;{\color{red}+}\;Q\,(1 - x^{v}_{ij}),
\]
which has the wrong sign of the big-$M$. With $+M$, every arc with
$x^{v}_{ij}=0$ would force $w^{v}_{j} \le w^{v}_{i} + q_{j} + \dots - Q$
(re-arranging), which is structurally infeasible whenever $Q$ is large.
The standard relaxation is $-M(1-x^{v}_{ij})$, which deactivates the
inequality whenever the arc is unused. We implement this minus sign.

\paragraph{Demand-sign convention.}
The paper's \S5.2 says ``demand $q_i$ is randomly drawn from $U(0, 2.27)$
with probability 0.86 \dots and the pickup demand is $-q_{n+i}$.'' Taking
this literally gives $q^{+} < 0$ at pickups and $q^{-} > 0$ at deliveries,
which contradicts (28) (truck load should \emph{rise} on a pickup). We
adopt the sign convention required by (28) and described in our
\texttt{instance\_generator.py}: $q_i > 0$ at pickup, $q_i < 0$ at delivery.

\paragraph{(26) Pickup-before-launch: depot launches.}
Constraint~(26) reads $u^{v}_{i} - u^{v}_{j-n} \ge 1 - M(1{-}y^{v}_{ijk})$
for $j\in\Phi^{-}$, $i\in C$. Two issues arise. First, $i$ runs over $C$
only, so $i=0$ is implicitly excluded; nothing in the paper prevents
launching a delivery sortie directly from the depot, which would let the
drone deliver before the truck has performed the corresponding pickup.
Second, when $i \in C$ but the corresponding pickup $j-n$ has not been
visited yet, $u^{v}_{j-n}$ is undefined relative to $u^{v}_{i}$. We
strengthen (26) by setting $y^{v}_{0jk}=0$ for every $j\in\Phi^{-}$,
$k\in\Nplus$. With one truck this rules out the only physically
non-implementable depot-launch case, and matches the paper's heuristic
(Algorithm~4) which only considers post-pickup launches.

\paragraph{(10) is subsumed by (9).}
Constraint~(10), $\sum_{i,j} y^{v}_{0jk} \le \sum_{l\ne k} x^{v}_{lk}$,
demands the rendezvous $k$ be on the truck route when the launch is from
the depot. Constraint~(9) already requires that, with the convention
$\mathbb{1}\{i\in\text{route}\}=1$ for $i=0$. We therefore implement only
(9) and document this elision.

\paragraph{(35) is omitted.}
Constraint~(35), $p^{v}_{0,j}=1$ for all $j\in C$, requires precedence
variables on the depot. We omit it because all our other depot-linked
constraints (the MTZ inequality, time/position bounds and (3)--(4)) place
the depot first implicitly. Including (35) would require tracking
$p^{v}_{0j}$ throughout, which inflates the variable count by $|C|$
without changing the feasible set.

\paragraph{$P$ pre-filtering.}
We pre-filter the universe of sorties $P = \{(i,j,k): i\ne 0,\;j \in
\Cprime,\;k\in\Nplus,\;i,j,k\text{ distinct}\}$ to those satisfying
$s_L + (d_{ij}+d_{jk})/s' + d^{\,\text{drone}}_{j} + s_R \le e$. Sorties
that fail this geometric check are forbidden by~(19) anyway, but
filtering them out shrinks the model by roughly an order of magnitude.

\section{Implementation}

The Python package is organised as
\begin{itemize}
\item \texttt{instance\_generator.py}: random-instance factory matching
\S5.2 of the paper (Table~1 defaults), plus a manual constructor used by
the verification suite.
\item \texttt{dapdp\_model.py}: the MILP, exposing
\texttt{solve\_dapdp(inst, time\_limit, ...)} that returns a
\texttt{Solution} dataclass.
\item \texttt{baseline.py}: solves the same MILP with $\Cprime=\emptyset$
to obtain $z_{\pdp}$ for SPDP computation.
\item \texttt{verification.py}: seven scenarios, each producing a
matplotlib figure plus an independently audited feasibility report.
\item \texttt{sensitivity.py}: parameter sweeps with CSV and PDF output.
\item \texttt{validation.py}: small-scale comparison with
$\overline{\mathit{SPDP}}$ values reported in the paper.
\end{itemize}

The solver is Gurobi 13.0.2 under the free restricted licence (size cap
$\sim$2000 vars/constraints). After the pre-filtering described above,
$n=4$ instances fit; $n=5$ does not.

\section{Verification}
\label{sec:verification}

Following the assignment guidelines, we devise instances whose feasible
solutions are restricted in a well-understood way and verify the solver
respects those restrictions. Each scenario is paired with an independent
auditor (\texttt{audit()} in \texttt{verification.py}) that re-derives
feasibility from the returned $x$ and $y$ values without using any
internal solver state.

\subsection{Scenarios}

\begin{description}[style=nextline,leftmargin=2em]
\item[V1, basic feasibility ($n{=}3$).] A randomly generated instance.
The auditor checks every constraint listed in
Table~\ref{tab:constraints}. Result: \emph{passed}, $z^{\star}=\$0.91$,
runtime $0.3$\,s. Figure~\ref{fig:V1}.

\item[V2, pickup-precedence (handcrafted).] All demands set to $5$\,kg
($>Q'$), so $\Cprime=\emptyset$ and the truck must serve every pair. The
auditor confirms each delivery follows its pickup in the truck sequence:
route $0\to1\to2\to3\to6\to5\to4\to7$. Figure~\ref{fig:V2}.

\item[V3, endurance.] The same instance solved with $e=60$\,min and
$e=5$\,min. With generous endurance the solver uses $2$ drone sorties
($z^{\star}=\$4.65$); with the $5$\,min limit no sortie is feasible
($z^{\star}=\$5.34$, $0$ sorties). Figure~\ref{fig:V3}. Both audits pass.

\item[V4, drone payload.] Three deliveries with demands $\{1, 50, 1\}$
kg. Only deliveries $4$ and $6$ are in $\Cprime$. The drone serves
exclusively delivery $6$; delivery $5$ (heavy) is on the truck route.
Figure~\ref{fig:V4}.

\item[V5, no sub-tours ($n{=}4$).] Verifies the truck's arc set forms a
single connected path $0\to\dots\to 2n{+}1$ with no detached cycles.
Figure~\ref{fig:V5}.

\item[V6, $z_{\dapdp}\le z_{\pdp}$.] Same instance ($n{=}4$, grid $10$,
seed $3$) solved twice. $z_{\pdp}=\$2.60$, $z_{\dapdp}=\$2.08$, a $20$\%
saving. Figure~\ref{fig:V6}.

\item[V7, no depot-launch deliveries.] After our extension of (26),
$y^{v}_{0jk}=0$ for $j\in\Phi^{-}$. The auditor confirms no drone sortie
in the optimum has $i=0$.
\end{description}

\begin{figure}[t]
\centering
\begin{subfigure}{0.46\textwidth}
\includegraphics[width=\linewidth]{figures/V1_basic.pdf}
\caption{V1 basic feasibility, $n{=}3$.}
\label{fig:V1}
\end{subfigure}\hfill
\begin{subfigure}{0.46\textwidth}
\includegraphics[width=\linewidth]{figures/V2_pickup_precedence.pdf}
\caption{V2 pickup-before-delivery (drone disabled).}
\label{fig:V2}
\end{subfigure}
\caption{Verification scenarios V1 and V2.}
\end{figure}

\begin{figure}[t]
\centering
\begin{subfigure}{0.46\textwidth}
\includegraphics[width=\linewidth]{figures/V3_endurance_60min.pdf}
\caption{$e=60$\,min: $2$ sorties, $z^{\star}=\$4.65$.}
\end{subfigure}\hfill
\begin{subfigure}{0.46\textwidth}
\includegraphics[width=\linewidth]{figures/V3_endurance_05min.pdf}
\caption{$e=5$\,min: $0$ sorties, $z^{\star}=\$5.34$.}
\end{subfigure}
\caption{V3 endurance test on the same instance.}
\label{fig:V3}
\end{figure}

\begin{figure}[t]
\centering
\begin{subfigure}{0.46\textwidth}
\includegraphics[width=\linewidth]{figures/V4_drone_payload.pdf}
\caption{V4: heavy delivery (5) excluded from $\Cprime$.}
\label{fig:V4}
\end{subfigure}\hfill
\begin{subfigure}{0.46\textwidth}
\includegraphics[width=\linewidth]{figures/V5_no_subtours.pdf}
\caption{V5 sub-tour elimination ($n{=}4$).}
\label{fig:V5}
\end{subfigure}
\caption{Verification scenarios V4 and V5.}
\end{figure}

\begin{figure}[t]
\centering
\begin{subfigure}{0.46\textwidth}
\includegraphics[width=\linewidth]{figures/V6_pdp.pdf}
\caption{Truck-only \pdp{}, $z=\$2.60$.}
\end{subfigure}\hfill
\begin{subfigure}{0.46\textwidth}
\includegraphics[width=\linewidth]{figures/V6_dapdp.pdf}
\caption{\dapdp{} optimum, $z=\$2.08$ ($20$\% saving).}
\end{subfigure}
\caption{V6: same instance, both formulations.}
\label{fig:V6}
\end{figure}

\subsection{Audit summary}
\begin{table}[h]
\centering
\caption{Verification suite results.}
\label{tab:verification}
\small
\begin{tabular}{@{}llrrl@{}}
\toprule
ID & description                       & $z^\star$ (\$) & runtime (s) & verdict \\
\midrule
V1 & basic, $n{=}3$                    & $0.912$ & $0.3$ & passed \\
V2 & pickup-precedence (truck only)    & $0.914$ & $0.5$ & passed \\
V3a & endurance 60 min                 & $4.652$ & $0.6$ & passed (2 sorties) \\
V3b & endurance 5 min                  & $5.341$ & $0.6$ & passed (0 sorties) \\
V4 & drone payload                     & $0.792$ & $0.4$ & passed (drone serves only $\{6\}$) \\
V5 & no sub-tours, $n{=}4$             & $3.847$ & $1.4$ & passed (single path) \\
V6 & DAPDP $\le$ PDP, $n{=}4$          & $2.078$ vs $2.603$ & $1.8$ & passed (20\% saving) \\
V7 & no depot-launch deliveries        & $0.664$ & $0.4$ & passed (no $y_{0jk}{=}1$) \\
\bottomrule
\end{tabular}
\end{table}

\section{Sensitivity analysis}
\label{sec:sensitivity}

We sweep three parameters that the paper highlights as most influential. The
ensemble is fixed at $18$ instances ($n\in\{3,4\}$, grid sizes
$\{5,10,20\}$\,miles, seeds $\{1,2,3\}$); $z_{\pdp}$ is computed once per
instance and cached, so each sweep value is one DAPDP solve per instance.

\subsection{Endurance}

Figure~\ref{fig:sens_endurance} shows mean SPDP rising monotonically with
endurance and saturating around $30$\,min, the paper's nominal value.
This matches the paper's Table~3 / Figure~9: when the drone can fly long
enough to serve more eligible deliveries, savings approach a ceiling
imposed by the drone's payload and route geometry.

\begin{figure}[h]
\centering
\includegraphics[width=0.6\textwidth]{figures/endurance.pdf}
\caption{Sensitivity of SPDP and number of sorties to drone endurance.}
\label{fig:sens_endurance}
\end{figure}

\subsection{Drone speed}

Figure~\ref{fig:sens_speed}: faster drones save more, but the relationship
flattens above $50$\,mph because most of the savings come from the
\emph{routing detour avoided} on the truck side rather than from the drone
itself. This too matches the paper.

\begin{figure}[h]
\centering
\includegraphics[width=0.6\textwidth]{figures/drone_speed.pdf}
\caption{Sensitivity of SPDP and number of sorties to drone speed.}
\label{fig:sens_speed}
\end{figure}

\subsection{Drone cost factor $\alpha$}

Figure~\ref{fig:sens_alpha}: SPDP decreases monotonically with $\alpha$.
At $\alpha=1$ no cost incentive remains and the model recovers $z_{\pdp}$
exactly (mean SPDP $\to 0$, mean sorties $\to 0$). At $\alpha=0.05$ the
mean saving is $\sim$21\%. The paper does not perform this sensitivity
explicitly; it is included because the assignment guidelines require
sensitivity to a cost parameter.

\begin{figure}[h]
\centering
\includegraphics[width=0.6\textwidth]{figures/alpha.pdf}
\caption{Sensitivity of SPDP and number of sorties to the drone cost
factor $\alpha = c'/c$.}
\label{fig:sens_alpha}
\end{figure}

\section{Validation}
\label{sec:validation}

A small-scale reproduction of the paper's Table~A.5 is given in
Table~\ref{tab:validation}. We solve $5$ random seeds per $(n, \text{grid})$
pair for $n \in \{3,4\}$ and grid in $\{5,10,20,30\}$\,miles. The free
licence prevents us from running $n=5$. Mean SPDP for $n\in\{3,4\}$ on
small/medium grids is $\sim 18$--$20$\%, in the same ballpark as the paper's
mean. The sharp drop at grid $30$\,mi is expected: the drone's $30$-min
endurance corresponds to $25$\,mi of round-trip flight, so most sorties
become geometrically infeasible at that scale.

\begin{table}[h]
\centering
\caption{Mean SPDP and runtime over $5$ random seeds per cell.}
\label{tab:validation}
\small
\begin{tabular}{@{}rrrrr@{}}
\toprule
$n$ & grid (mi) & mean SPDP (\%) & std (\%) & mean runtime (s) \\
\midrule
3 &  5 & 18.5 & 6.2 & 0.14 \\
3 & 10 & 18.5 & 6.2 & 0.17 \\
3 & 20 & 16.1 & 6.2 & 0.19 \\
3 & 30 &  7.2 & 4.9 & 0.09 \\
4 &  5 & 20.1 & 12.4 & 0.64 \\
4 & 10 & 20.1 & 12.4 & 0.78 \\
4 & 20 & 18.6 & 12.4 & 0.65 \\
4 & 30 &  4.6 & 3.0 & 0.30 \\
\bottomrule
\end{tabular}
\end{table}

\section{Conclusions}

A faithful Python/Gurobi reproduction of the Mulumba--Diabat \dapdp{} MILP
is presented and verified. The seven verification scenarios all pass
their independent audits; SPDP increases with drone endurance and speed
and decreases with $\alpha$, in line with the paper's findings; mean
savings of $\sim$\,18--20\% are recovered on small-scale instances. Three
typographic ambiguities of the paper are documented and resolved.

The main practical limitation of this implementation is the size cap
imposed by the free Gurobi licence, which restricts us to $n\le 4$. The
underlying model scales as expected; an academic licence would lift this
limit immediately. Future work would be to wire in the paper's ALNS
heuristic (Section~4), which would enable validation against the larger
test instances of Section~5.

\begin{thebibliography}{1}
\bibitem{mulumba2024}
T.~Mulumba and A.~Diabat. \emph{Optimization of the drone-assisted pickup
and delivery problem}. Transportation Research Part E, 181:103377, 2024.
\end{thebibliography}

\end{document}
